% problem_id: S046-thm:trunc-in-truncated-sigma
% source_id: S046 (HoTT Book)
% source_file: hlevels.tex
% statement_locator: hlevels.tex lines 629--634
% proof_locator: hlevels.tex lines 636--649
% env: thm
% pairing: tight (verified)
% status: complete (statement + author proof verbatim + reason + locators)
%
\begin{thm}\label{thm:trunc-in-truncated-sigma}
Let $P:A\to\type$ be a family of types. Then there is an equivalence
\begin{equation*}
\eqv{\Trunc n{\sm{x:A}\trunc n{P(x)}}}{\Trunc n{\sm{x:A}P(x)}}.
\end{equation*}
\end{thm}

\begin{proof}
We use the induction principle of $n$-truncation several times to construct
functions
\begin{align*}
\varphi & : \Trunc n{\sm{x:A}\trunc n{P(x)}}\to\Trunc n{\sm{x:A}P(x)}\\
\psi & : \Trunc n{\sm{x:A}P(x)}\to \Trunc n{\sm{x:A} \trunc n{P(x)}}
\end{align*}
and homotopies $H:\varphi\circ\psi\htpy \idfunc$ and $K:\psi\circ\varphi\htpy
\idfunc$ exhibiting them as quasi-inverses.
We define $\varphi$ by setting $\varphi(\tproj n{\pairr{x,\tproj nu}})\defeq\tproj n{\pairr{x,u}}$.
We define $\psi$ by setting $\psi(\tproj n{\pairr{x,u}})\defeq\tproj n{\pairr{x,\tproj nu}}$.
Then we define $H(\tproj n{\pairr{x,u}})\defeq \refl{\tproj n{\pairr{x,u}}}$ and
$K(\tproj n{\pairr{x,\tproj nu}})\defeq \refl{\tproj n{\pairr{x,\tproj nu}}}$.
\end{proof}
